\label{ch:gauss}

\section{Cofinalidad entre cilindros HMT y cilindros de fracciones continuas}

Una secci\'on HMT irracional determina una familia anidada de cilindros compactos. Su evaluaci\'on arquimediana es un punto \(x\in(0,1)\setminus\mathbb Q\). La fracci\'on continua de \(x\) determina cilindros cl\'asicos

\[
I(a_1,\ldots,a_n).
\]

\begin{proposition}[Cofinalidad]
Para todo cilindro finito \(I(a_1,\ldots,a_n)\) que contiene la evaluaci\'on \(x\), existe profundidad HMT \(m(n)\) tal que el cilindro HMT de nivel \(m\ge m(n)\) est\'a contenido en \(I(a_1,\ldots,a_n)\).
\end{proposition}

\begin{proof}
Fijado \(n\), el cilindro de fracci\'on continua \(I(a_1,\ldots,a_n)\) es un entorno abierto de la irracional \(x\). Existe entonces \(\varepsilon_n>0\) con \((x-\varepsilon_n,x+\varepsilon_n)\subset I(a_1,\ldots,a_n)\). Los cilindros HMT \(C_m\) contienen \(x\), son anidados y satisfacen \(\operatorname{diam}C_m\to0\). Elija \(m(n)\) con \(\operatorname{diam}C_{m(n)}<\varepsilon_n\). Como \(x\in C_{m(n)}\), todo punto de ese cilindro dista de \(x\) menos de \(\varepsilon_n\), luego \(C_{m(n)}\subset I(a_1,\ldots,a_n)\). El anidamiento da la inclusi\'on para todo nivel posterior.
\end{proof}

La demostraci\'on utiliza compacidad, anidamiento y unicidad de la
evaluaci\'on. El resultado establece el transporte de aproximaciones, pero
no implica la conjugaci\'on del desplazamiento de memoria TPK con el mapa de
Gauss sobre todas las fibras profundas.

La genealog\'ia completa de esta realizaci\'on es
\[
\APP\to\TRIT\to\TPK
\to (C_m,\rho_{m+1,m})_{m\ge0}
\to x=\bigcap_m C_m
\to \bigl(I(a_1,\ldots,a_n)\bigr)_{n\ge1}.
\]
Los cilindros HMT conservan ruta, fase y acarreo; el mapa terminal proyecta
estas coordenadas sobre los prefijos de fracci\'on continua. La cofinalidad
demuestra que dicha proyecci\'on converge al mismo punto. Un entrelazador
din\'amico exigir\'ia adem\'as un levantamiento \(\widetilde T_G\) con
cuadrados conmutativos sobre las fibras de valoraci\'on, propiedad que no se
deduce de la cofinalidad.

\section{Medida de Gauss y observables de Khinchin y L\'evy}

El mapa de Gauss \cite{Khinchin,IosifescuKraaikamp}

\[
T_G(x)=\frac1x-\left\lfloor\frac1x\right\rfloor
\]

preserva

\begin{equation}
d\mu_G(x)=\frac{dx}{\log2(1+x)}.
\label{eq:gauss-measure}
\end{equation}

La constante de Khinchin se define mediante

\begin{equation}
\log K=\int_0^1\log a_1(x)\,d\mu_G(x).
\label{eq:khinchin}
\end{equation}

La distribuci\'on del primer d\'igito se calcula directamente:
\begin{equation}
\mu_G(a_1=k)
=\frac1{\log2}\int_{1/(k+1)}^{1/k}\frac{dx}{1+x}
=\log_2\!\left(\frac{(k+1)^2}{k(k+2)}\right).
\label{eq:gauss-digit-law}
\end{equation}
Por el teorema erg\'odico,
\begin{equation}
K=\prod_{k\ge1}
k^{\,\log_2((k+1)^2/[k(k+2)])}
=2.685452001065306445309714835\ldots
\label{eq:khinchin-product}
\end{equation}
para \(\mu_G\)-casi toda secci\'on. El producto no se aplica a cada punto:
la secci\'on \(\varphi^{-1}\), por ejemplo, tiene todos sus d\'igitos iguales
a uno y es un criterio de refutación inmediato de una lectura punto a punto. Una
certificaci\'on de \cref{eq:khinchin-product} trunca la suma de logaritmos y
acota la cola mediante la estimaci\'on
\[
\mu_G(a_1=k)=O(k^{-2}),
\qquad
\sum_{k>N}\mu_G(a_1=k)\log k
=O\!\left(\frac{\log N+1}{N}\right).
\]

La constante de L\'evy es

\begin{equation}
L=-\int_0^1\log x\,d\mu_G(x)
=\frac{\pi^2}{12\log2}.
\label{eq:levy}
\end{equation}

En efecto,

\[
\frac1{1+x}=\sum_{n\ge0}(-1)^nx^n,\qquad
\int_0^1(-\log x)x^n\,dx=\frac1{(n+1)^2},
\]

por lo que la integral es \(\eta(2)/\log2=\pi^2/(12\log2)\).

La constante de Khinchin determina la media logar\'itmica de los cocientes
parciales, mientras que la constante de L\'evy determina la tasa exponencial
de crecimiento de los denominadores. Ambas son observables distintos de un
mismo sistema medido.

\section{Entrop\'ia m\'etrica, exponente de Lyapunov y constante de Lochs}

Como

\[
\log|T_G'(x)|=-2\log x,
\]

el exponente de Lyapunov y la entrop\'ia m\'etrica son

\begin{equation}
\chi_G=h_G=2L=\frac{\pi^2}{6\log2}.
\label{eq:gauss-entropy}
\end{equation}

La raz\'on de Lochs entre fracciones continuas y una expansi\'on en base \(b\) es

\begin{equation}
\Lambda_b=\frac{\log b}{h_G}
=\frac{6\log2\log b}{\pi^2}.
\label{eq:lochs}
\end{equation}

En particular,

\begin{align*}
\Lambda_{10}&=0.9702701143920339257402560192\ldots,\\
\Lambda_{729}&=\frac{36\log2\log3}{\pi^2}
=2.777618966374305777705782976\ldots.
\end{align*}

La base \(729=3^6\) es relevante para el ambiente \(\mathbb F_3^6\), pero \cref{eq:lochs} sigue siendo un teorema casi seguro de \(\mu_G\); no una igualdad punto a punto para toda ruta HMT.

El valor \(\Lambda_{729}\) representa la tasa asint\'otica t\'ipica de
cocientes parciales obtenidos por bloque completo de seis trits. No identifica el
ambiente de \(729\) palabras con los otros objetos de cardinal \(729\) del
corpus. La igualdad es una tasa de informaci\'on entre particiones, no una
bijecci\'on entre estados.

\section{Operador de transferencia de Gauss--Kuzmin--Wirsing}

El operador de Gauss--Kuzmin--Wirsing act\'ua por

\begin{equation}
(\mathcal L f)(x)=\sum_{n\ge1}\frac1{(x+n)^2}
f\!\left(\frac1{x+n}\right).
\label{eq:gkw-operator}
\end{equation}

Su autovalor dominante corresponde a la densidad invariante. La constante GKW es el m\'odulo del autovalor subdominante sobre el subespacio de media cero. Un esquema HMT compatible usa particiones de \(9^m\) cilindros, matrices de transferencia con intervalos y una cota de cola.

Para promover el c\'alculo se deben certificar: aislamiento del autovalor, invariancia del subespacio, error de truncaci\'on, convergencia de proyectores espectrales y ausencia de otro autovalor de igual m\'odulo. El decimal conocido no puede introducirse como centro del intervalo inicial.

Una deformaci\'on conjunta contiene Khinchin y L\'evy en un solo objeto:
\begin{equation}
(\mathcal L_{s,t}f)(x)
=\sum_{n\ge1}n^t(n+x)^{-2s}
 f\!\left(\frac1{n+x}\right),
\qquad
P(s,t)=\log r(\mathcal L_{s,t}).
\label{eq:gauss-pressure-operator}
\end{equation}
En el punto \((s,t)=(1,0)\), la perturbaci\'on espectral da
\begin{equation}
\partial_tP(1,0)=\log K,
\qquad
-\partial_sP(1,0)=\chi_G=2L.
\label{eq:pressure-derivatives}
\end{equation}
La Hessiana de \(P\) determina las varianzas y covarianzas asint\'oticas de
\(\log a_1\) y \(-2\log x\). Por tanto, las constantes de Khinchin y
L\'evy son dos derivadas transversales de la misma funci\'on de presi\'on.

El reconocimiento num\'erico del autovalor subdominante es
\[
\lambda_2\simeq-0.3036630028987326586,\qquad
\kappa_{\rm GKW}=|\lambda_2|,\qquad
-\log\kappa_{\rm GKW}\simeq1.19183673556055.
\]
En este volumen, dichos decimales constituyen valores de referencia para el
contraste y no un certificado espectral.
La prueba HMT propuesta usa Galerkin o Ulam sobre los cilindros cofinales de
profundidad \(9^m\), una desigualdad de Lasota--Yorke, aritm\'etica de
intervalos para aislar \(\lambda_2\) y una cota rigurosa de la cola. Un
intervalo que contenga dos autovalores del mismo m\'odulo refuta la
promoci\'on, aunque no afecte a Khinchin, L\'evy o Lochs.

